% ============================================================
% Scenario A: BN artifact confirmed (Detach+LN CAR ≈ Vanilla ~0.780)
% Probability: ~40%
% Key implication: Detach degradation was normalization artifact, not gradient flow evidence
% ============================================================

% REPLACES: "We hypothesize..." paragraph in §6.2 (discussion.tex L30)
% ALSO UPDATES: preceding "Gradient analysis" paragraph

\begin{comment} % --- SCENARIO A START ---

A discriminating test replaces the auxiliary head's BatchNorm2d layers with LayerNorm, which normalizes per-sample and does not maintain cross-batch running statistics, while keeping gradient detachment (\texttt{stop\_gradient}) unchanged. Detach+LN (3 seeds) achieves CAR of [DETACH_LN_CAR] ($\pm$ [DETACH_LN_STD]), recovering to within [DELTA]pp of Vanilla (0.780) and statistically indistinguishable from the no-auxiliary baseline. This confirms that the Detach degradation is a normalization artifact: BatchNorm's running statistics create forward-pass interference that the encoder cannot compensate for without auxiliary gradients. When this interference is removed via LayerNorm, the encoder performs comparably to Vanilla even without auxiliary gradient flow.

This result refines the gradient-flow finding. Gradient flow is not inherently beneficial for collision avoidance---it is necessary specifically to compensate for the distributional perturbation that BatchNorm layers introduce in the auxiliary head. The practical implication is twofold: (1)~if practitioners use auxiliary heads with BatchNorm during VLA fine-tuning, gradient flow through the encoder is mandatory to avoid degradation; (2)~if the auxiliary head uses per-sample normalization (LayerNorm, GroupNorm), gradient detachment is safe and does not harm collision avoidance. This normalization-dependent finding explains why all gradient-flow conditions (Full OCC, Random OCC, Depth) achieve comparable CAR regardless of label content: the gradient signal compensates for the same BN artifact in all three cases.

\end{comment} % --- SCENARIO A END ---


% ============================================================
% Scenario B: Gradient flow is genuine (Detach+LN CAR ≈ Detach ~0.746)
% Probability: ~45%
% Key implication: Gradient flow itself protects collision avoidance, independent of normalization
% ============================================================

% REPLACES: same paragraph in §6.2

\begin{comment} % --- SCENARIO B START ---

To test whether BatchNorm interference accounts for the Detach degradation, we replace the auxiliary head's BatchNorm2d layers with LayerNorm---which normalizes per-sample and does not maintain cross-batch running statistics---while keeping gradient detachment unchanged. Detach+LN (3 seeds) achieves CAR of [DETACH_LN_CAR] ($\pm$ [DETACH_LN_STD]), comparable to Detach OCC (0.746) and well below Vanilla (0.780). Replacing BatchNorm with LayerNorm does not restore collision avoidance, ruling out normalization interference as the primary mechanism. The gradient-flow effect is genuine: blocking encoder gradients from the auxiliary loss degrades collision avoidance regardless of the normalization scheme in the auxiliary head.

This result strengthens the central finding. Gradient flow through the vision encoder is a necessary condition for auxiliary heads to avoid degrading collision avoidance, and this necessity is independent of architectural details in the auxiliary pathway. The three candidate mechanisms proposed in Section~\ref{sec:why_gradient}---anti-overfitting regularization, optimization landscape smoothing, and gradient diversity---all remain viable explanations, since each predicts that gradient flow matters independently of normalization. The Detach+LN control eliminates the confound that BatchNorm artifacts could have driven the Detach--Vanilla gap, establishing gradient flow absence as the operative factor.

\end{comment} % --- SCENARIO B END ---


% ============================================================
% Scenario C: Mixed (Detach+LN CAR between Detach and Vanilla, e.g., ~0.760)
% Probability: ~15%
% Key implication: Both BN interference and gradient flow absence contribute
% ============================================================

% REPLACES: same paragraph in §6.2

\begin{comment} % --- SCENARIO C START ---

To decompose the Detach degradation into normalization and gradient-flow components, we replace the auxiliary head's BatchNorm2d layers with LayerNorm while keeping gradient detachment unchanged. Detach+LN (3 seeds) achieves CAR of [DETACH_LN_CAR] ($\pm$ [DETACH_LN_STD]), falling between Detach OCC (0.746) and Vanilla (0.780). Removing BatchNorm interference recovers approximately [BN_FRACTION]\% of the Detach--Vanilla gap, indicating that both normalization artifacts and gradient-flow absence contribute to the degradation. The remaining [GF_FRACTION]\% gap between Detach+LN and Vanilla reflects the genuine cost of blocking auxiliary gradients from the encoder, independent of normalization design.

This partial recovery yields a more nuanced account than either pure mechanism alone would predict. BatchNorm's running statistics introduce forward-pass interference that accounts for roughly [BN_PP]pp of the 3.4pp Detach deficit, while gradient-flow absence accounts for the remaining [GF_PP]pp. The practical implication is that normalization choice in auxiliary heads matters---using LayerNorm instead of BatchNorm mitigates but does not eliminate the cost of gradient detachment. For maximal safety preservation, both design choices are needed: per-sample normalization to avoid batch-statistics contamination, and gradient flow to maintain the regularization benefit of multi-task training.

\end{comment} % --- SCENARIO C END ---


% ============================================================
% Corresponding Claim 1 variants for Abstract/Introduction
% ============================================================

% SCENARIO A - C1:
% "When using BatchNorm in auxiliary heads, gradient flow is necessary to
%  compensate for running statistics interference; with per-sample normalization,
%  gradient detachment is safe."

% SCENARIO B - C1 (strengthened):
% "Gradient flow through the vision encoder is a necessary condition for
%  auxiliary heads to avoid degrading collision avoidance, independent of
%  normalization architecture."

% SCENARIO C - C1:
% "Both gradient flow and normalization design in auxiliary heads contribute
%  to collision avoidance preservation; gradient flow is necessary but
%  normalization choice modulates the severity of detachment."
